% Options for packages loaded elsewhere
\PassOptionsToPackage{unicode}{hyperref}
\PassOptionsToPackage{hyphens}{url}
\documentclass[
]{article}
\usepackage{xcolor}
\usepackage[margin=1in]{geometry}
\usepackage{amsmath,amssymb}
\setcounter{secnumdepth}{-\maxdimen} % remove section numbering
\usepackage{iftex}
\ifPDFTeX
  \usepackage[T1]{fontenc}
  \usepackage[utf8]{inputenc}
  \usepackage{textcomp} % provide euro and other symbols
\else % if luatex or xetex
  \usepackage{unicode-math} % this also loads fontspec
  \defaultfontfeatures{Scale=MatchLowercase}
  \defaultfontfeatures[\rmfamily]{Ligatures=TeX,Scale=1}
\fi
\usepackage{lmodern}
\ifPDFTeX\else
  % xetex/luatex font selection
\fi
% Use upquote if available, for straight quotes in verbatim environments
\IfFileExists{upquote.sty}{\usepackage{upquote}}{}
\IfFileExists{microtype.sty}{% use microtype if available
  \usepackage[]{microtype}
  \UseMicrotypeSet[protrusion]{basicmath} % disable protrusion for tt fonts
}{}
\makeatletter
\@ifundefined{KOMAClassName}{% if non-KOMA class
  \IfFileExists{parskip.sty}{%
    \usepackage{parskip}
  }{% else
    \setlength{\parindent}{0pt}
    \setlength{\parskip}{6pt plus 2pt minus 1pt}}
}{% if KOMA class
  \KOMAoptions{parskip=half}}
\makeatother
\usepackage{color}
\usepackage{fancyvrb}
\newcommand{\VerbBar}{|}
\newcommand{\VERB}{\Verb[commandchars=\\\{\}]}
\DefineVerbatimEnvironment{Highlighting}{Verbatim}{commandchars=\\\{\}}
% Add ',fontsize=\small' for more characters per line
\usepackage{framed}
\definecolor{shadecolor}{RGB}{248,248,248}
\newenvironment{Shaded}{\begin{snugshade}}{\end{snugshade}}
\newcommand{\AlertTok}[1]{\textcolor[rgb]{0.94,0.16,0.16}{#1}}
\newcommand{\AnnotationTok}[1]{\textcolor[rgb]{0.56,0.35,0.01}{\textbf{\textit{#1}}}}
\newcommand{\AttributeTok}[1]{\textcolor[rgb]{0.13,0.29,0.53}{#1}}
\newcommand{\BaseNTok}[1]{\textcolor[rgb]{0.00,0.00,0.81}{#1}}
\newcommand{\BuiltInTok}[1]{#1}
\newcommand{\CharTok}[1]{\textcolor[rgb]{0.31,0.60,0.02}{#1}}
\newcommand{\CommentTok}[1]{\textcolor[rgb]{0.56,0.35,0.01}{\textit{#1}}}
\newcommand{\CommentVarTok}[1]{\textcolor[rgb]{0.56,0.35,0.01}{\textbf{\textit{#1}}}}
\newcommand{\ConstantTok}[1]{\textcolor[rgb]{0.56,0.35,0.01}{#1}}
\newcommand{\ControlFlowTok}[1]{\textcolor[rgb]{0.13,0.29,0.53}{\textbf{#1}}}
\newcommand{\DataTypeTok}[1]{\textcolor[rgb]{0.13,0.29,0.53}{#1}}
\newcommand{\DecValTok}[1]{\textcolor[rgb]{0.00,0.00,0.81}{#1}}
\newcommand{\DocumentationTok}[1]{\textcolor[rgb]{0.56,0.35,0.01}{\textbf{\textit{#1}}}}
\newcommand{\ErrorTok}[1]{\textcolor[rgb]{0.64,0.00,0.00}{\textbf{#1}}}
\newcommand{\ExtensionTok}[1]{#1}
\newcommand{\FloatTok}[1]{\textcolor[rgb]{0.00,0.00,0.81}{#1}}
\newcommand{\FunctionTok}[1]{\textcolor[rgb]{0.13,0.29,0.53}{\textbf{#1}}}
\newcommand{\ImportTok}[1]{#1}
\newcommand{\InformationTok}[1]{\textcolor[rgb]{0.56,0.35,0.01}{\textbf{\textit{#1}}}}
\newcommand{\KeywordTok}[1]{\textcolor[rgb]{0.13,0.29,0.53}{\textbf{#1}}}
\newcommand{\NormalTok}[1]{#1}
\newcommand{\OperatorTok}[1]{\textcolor[rgb]{0.81,0.36,0.00}{\textbf{#1}}}
\newcommand{\OtherTok}[1]{\textcolor[rgb]{0.56,0.35,0.01}{#1}}
\newcommand{\PreprocessorTok}[1]{\textcolor[rgb]{0.56,0.35,0.01}{\textit{#1}}}
\newcommand{\RegionMarkerTok}[1]{#1}
\newcommand{\SpecialCharTok}[1]{\textcolor[rgb]{0.81,0.36,0.00}{\textbf{#1}}}
\newcommand{\SpecialStringTok}[1]{\textcolor[rgb]{0.31,0.60,0.02}{#1}}
\newcommand{\StringTok}[1]{\textcolor[rgb]{0.31,0.60,0.02}{#1}}
\newcommand{\VariableTok}[1]{\textcolor[rgb]{0.00,0.00,0.00}{#1}}
\newcommand{\VerbatimStringTok}[1]{\textcolor[rgb]{0.31,0.60,0.02}{#1}}
\newcommand{\WarningTok}[1]{\textcolor[rgb]{0.56,0.35,0.01}{\textbf{\textit{#1}}}}
\usepackage{graphicx}
\makeatletter
\newsavebox\pandoc@box
\newcommand*\pandocbounded[1]{% scales image to fit in text height/width
  \sbox\pandoc@box{#1}%
  \Gscale@div\@tempa{\textheight}{\dimexpr\ht\pandoc@box+\dp\pandoc@box\relax}%
  \Gscale@div\@tempb{\linewidth}{\wd\pandoc@box}%
  \ifdim\@tempb\p@<\@tempa\p@\let\@tempa\@tempb\fi% select the smaller of both
  \ifdim\@tempa\p@<\p@\scalebox{\@tempa}{\usebox\pandoc@box}%
  \else\usebox{\pandoc@box}%
  \fi%
}
% Set default figure placement to htbp
\def\fps@figure{htbp}
\makeatother
\setlength{\emergencystretch}{3em} % prevent overfull lines
\providecommand{\tightlist}{%
  \setlength{\itemsep}{0pt}\setlength{\parskip}{0pt}}
\usepackage{bookmark}
\IfFileExists{xurl.sty}{\usepackage{xurl}}{} % add URL line breaks if available
\urlstyle{same}
% fallback for those not using the hyperref driver hyperxmp:
\makeatletter
\@ifundefined{xmpquote}{\newcommand{\xmpquote}[1]{#1}}{}
\makeatother
\hypersetup{
  pdftitle={Midterm Exam},
  pdfauthor={José Manuel Murcia},
  hidelinks,
  pdfcreator={LaTeX via pandoc}}

\title{Midterm Exam}
\author{José Manuel Murcia}
\date{2026-10-22}

\begin{document}
\maketitle

\section{Hola, soy Manuel Murcia}\label{hola-soy-manuel-murcia}

\subsubsection{Estudiante de la UPR
Mayagüez}\label{estudiante-de-la-upr-mayaguxfcez}

\paragraph{Maestría en Biología}\label{maestruxeda-en-biologuxeda}

\subsubsection{Examen 1: Jueves 22 de Octubre, 3:00 pm,
finalización:}\label{examen-1-jueves-22-de-octubre-300-pm-finalizaciuxf3n}

Aquí se muestran mis ejercicios en el aprendizaje de R. Examen1:

\begin{Shaded}
\begin{Highlighting}[]
\CommentTok{\# install.packages("tidyverse")}
\CommentTok{\# installed.packages("dplyr")}
\CommentTok{\# install.packages("aod")}

\FunctionTok{library}\NormalTok{(tidyverse)}
\end{Highlighting}
\end{Shaded}

\begin{verbatim}
-- Attaching core tidyverse packages ------------------------ tidyverse 2.0.0 --
v dplyr     1.2.1     v readr     2.2.0
v forcats   1.0.1     v stringr   1.6.0
v ggplot2   4.0.3     v tibble    3.3.1
v lubridate 1.9.5     v tidyr     1.3.2
v purrr     1.2.2     
-- Conflicts ------------------------------------------ tidyverse_conflicts() --
x dplyr::filter() masks stats::filter()
x dplyr::lag()    masks stats::lag()
i Use the conflicted package (<http://conflicted.r-lib.org/>) to force all conflicts to become errors
\end{verbatim}

\begin{Shaded}
\begin{Highlighting}[]
\FunctionTok{library}\NormalTok{(dplyr)}
\FunctionTok{library}\NormalTok{(aod)}

\FunctionTok{data}\NormalTok{(}\StringTok{"msleep"}\NormalTok{)}
\FunctionTok{data}\NormalTok{(}\StringTok{"storms"}\NormalTok{)}
\FunctionTok{data}\NormalTok{(}\StringTok{"lizards"}\NormalTok{)}

\FunctionTok{head}\NormalTok{(msleep, }\DecValTok{6}\NormalTok{)}
\end{Highlighting}
\end{Shaded}

\begin{verbatim}
# A tibble: 6 x 11
  name    genus vore  order conservation sleep_total sleep_rem sleep_cycle awake
  <chr>   <chr> <chr> <chr> <chr>              <dbl>     <dbl>       <dbl> <dbl>
1 Cheetah Acin~ carni Carn~ lc                  12.1      NA        NA      11.9
2 Owl mo~ Aotus omni  Prim~ <NA>                17         1.8      NA       7  
3 Mounta~ Aplo~ herbi Rode~ nt                  14.4       2.4      NA       9.6
4 Greate~ Blar~ omni  Sori~ lc                  14.9       2.3       0.133   9.1
5 Cow     Bos   herbi Arti~ domesticated         4         0.7       0.667  20  
6 Three-~ Brad~ herbi Pilo~ <NA>                14.4       2.2       0.767   9.6
# i 2 more variables: brainwt <dbl>, bodywt <dbl>
\end{verbatim}

\begin{Shaded}
\begin{Highlighting}[]
\FunctionTok{head}\NormalTok{(storms, }\DecValTok{6}\NormalTok{)}
\end{Highlighting}
\end{Shaded}

\begin{verbatim}
# A tibble: 6 x 13
  name   year month   day  hour   lat  long status       category  wind pressure
  <chr> <dbl> <dbl> <int> <dbl> <dbl> <dbl> <fct>           <dbl> <int>    <int>
1 Amy    1975     6    27     0  27.5 -79   tropical de~       NA    25     1013
2 Amy    1975     6    27     6  28.5 -79   tropical de~       NA    25     1013
3 Amy    1975     6    27    12  29.5 -79   tropical de~       NA    25     1013
4 Amy    1975     6    27    18  30.5 -79   tropical de~       NA    25     1013
5 Amy    1975     6    28     0  31.5 -78.8 tropical de~       NA    25     1012
6 Amy    1975     6    28     6  32.4 -78.7 tropical de~       NA    25     1012
# i 2 more variables: tropicalstorm_force_diameter <int>,
#   hurricane_force_diameter <int>
\end{verbatim}

\begin{Shaded}
\begin{Highlighting}[]
\FunctionTok{head}\NormalTok{(lizards, }\DecValTok{6}\NormalTok{)}
\end{Highlighting}
\end{Shaded}

\begin{verbatim}
  Site Diameter Height    Time grahami opalinus
1  Sun   D <= 2  H < 5   Early      20        2
2  Sun   D <= 2  H < 5 Mid-day       8        1
3  Sun   D <= 2  H < 5    Late       4        4
4  Sun   D <= 2 H >= 5   Early      13        0
5  Sun   D <= 2 H >= 5 Mid-day       8        0
6  Sun   D <= 2 H >= 5    Late      12        0
\end{verbatim}

\paragraph{Q1. Usa filter() para seleccionar del dataset msleep solo los
animales que están despiertos al menos 12 horas al
día.}\label{q1.-usa-filter-para-seleccionar-del-dataset-msleep-solo-los-animales-que-estuxe1n-despiertos-al-menos-12-horas-al-duxeda.}

\begin{Shaded}
\begin{Highlighting}[]
\NormalTok{msleep }\SpecialCharTok{\%\textgreater{}\%}
  \FunctionTok{filter}\NormalTok{(awake }\SpecialCharTok{\textgreater{}=} \DecValTok{12}\NormalTok{)}
\end{Highlighting}
\end{Shaded}

\begin{verbatim}
## # A tibble: 52 x 11
##    name   genus vore  order conservation sleep_total sleep_rem sleep_cycle awake
##    <chr>  <chr> <chr> <chr> <chr>              <dbl>     <dbl>       <dbl> <dbl>
##  1 Cow    Bos   herbi Arti~ domesticated         4         0.7       0.667  20  
##  2 North~ Call~ carni Carn~ vu                   8.7       1.4       0.383  15.3
##  3 Vespe~ Calo~ <NA>  Rode~ <NA>                 7        NA        NA      17  
##  4 Dog    Canis carni Carn~ domesticated        10.1       2.9       0.333  13.9
##  5 Roe d~ Capr~ herbi Arti~ lc                   3        NA        NA      21  
##  6 Goat   Capri herbi Arti~ lc                   5.3       0.6      NA      18.7
##  7 Guine~ Cavis herbi Rode~ domesticated         9.4       0.8       0.217  14.6
##  8 Grivet Cerc~ omni  Prim~ lc                  10         0.7      NA      14  
##  9 Star-~ Cond~ omni  Sori~ lc                  10.3       2.2      NA      13.7
## 10 Afric~ Cric~ omni  Rode~ <NA>                 8.3       2        NA      15.7
## # i 42 more rows
## # i 2 more variables: brainwt <dbl>, bodywt <dbl>
\end{verbatim}

\paragraph{Q2. Usa mutate() para crear una nueva columna llamada
percent\_day\_awake que muestre el porcentaje del día que cada especie
permanece despierta. Usa select() con everything() al final para
comprobar que tus cálculos son
correctos.}\label{q2.-usa-mutate-para-crear-una-nueva-columna-llamada-percent_day_awake-que-muestre-el-porcentaje-del-duxeda-que-cada-especie-permanece-despierta.-usa-select-con-everything-al-final-para-comprobar-que-tus-cuxe1lculos-son-correctos.}

\begin{Shaded}
\begin{Highlighting}[]
\NormalTok{  msleep }\OtherTok{\textless{}{-}}\NormalTok{  msleep }\SpecialCharTok{\%\textgreater{}\%} 
  \FunctionTok{mutate}\NormalTok{(}\AttributeTok{percent\_day\_awake =}\NormalTok{ (awake }\SpecialCharTok{/} \DecValTok{24}\NormalTok{) }\SpecialCharTok{*} \DecValTok{100}\NormalTok{) }\SpecialCharTok{\%\textgreater{}\%} 
  \FunctionTok{select}\NormalTok{(percent\_day\_awake, awake, }\FunctionTok{everything}\NormalTok{())}
\NormalTok{msleep}
\end{Highlighting}
\end{Shaded}

\begin{verbatim}
## # A tibble: 83 x 12
##    percent_day_awake awake name       genus vore  order conservation sleep_total
##                <dbl> <dbl> <chr>      <chr> <chr> <chr> <chr>              <dbl>
##  1              49.6  11.9 Cheetah    Acin~ carni Carn~ lc                  12.1
##  2              29.2   7   Owl monkey Aotus omni  Prim~ <NA>                17  
##  3              40     9.6 Mountain ~ Aplo~ herbi Rode~ nt                  14.4
##  4              37.9   9.1 Greater s~ Blar~ omni  Sori~ lc                  14.9
##  5              83.3  20   Cow        Bos   herbi Arti~ domesticated         4  
##  6              40     9.6 Three-toe~ Brad~ herbi Pilo~ <NA>                14.4
##  7              63.8  15.3 Northern ~ Call~ carni Carn~ vu                   8.7
##  8              70.8  17   Vesper mo~ Calo~ <NA>  Rode~ <NA>                 7  
##  9              57.9  13.9 Dog        Canis carni Carn~ domesticated        10.1
## 10              87.5  21   Roe deer   Capr~ herbi Arti~ lc                   3  
## # i 73 more rows
## # i 4 more variables: sleep_rem <dbl>, sleep_cycle <dbl>, brainwt <dbl>,
## #   bodywt <dbl>
\end{verbatim}

\paragraph{Q3. Agrupa los datos por año y calcula el número de tormentas
que ocurrieron cada año. Sugerencia: group\_by(col) agrupa las filas por
tiempo, y summarise(num\_storms = n()) cuenta cuántas filas (tormentas)
hay por
año.}\label{q3.-agrupa-los-datos-por-auxf1o-y-calcula-el-nuxfamero-de-tormentas-que-ocurrieron-cada-auxf1o.-sugerencia-group_bycol-agrupa-las-filas-por-tiempo-y-summarisenum_storms-n-cuenta-cuuxe1ntas-filas-tormentas-hay-por-auxf1o.}

\begin{Shaded}
\begin{Highlighting}[]
\NormalTok{storms }\SpecialCharTok{\%\textgreater{}\%}
  \FunctionTok{group\_by}\NormalTok{(year) }\SpecialCharTok{\%\textgreater{}\%}
  \FunctionTok{summarise}\NormalTok{(}
    \AttributeTok{num\_storms =} \FunctionTok{n}\NormalTok{()}
\NormalTok{  )}
\end{Highlighting}
\end{Shaded}

\begin{verbatim}
## # A tibble: 50 x 2
##     year num_storms
##    <dbl>      <int>
##  1  1975        238
##  2  1976        126
##  3  1977         92
##  4  1978        152
##  5  1979        324
##  6  1980        335
##  7  1981        311
##  8  1982        111
##  9  1983         88
## 10  1984        342
## # i 40 more rows
\end{verbatim}

\paragraph{Q4 Filtra el dataset para quedarte solo con tormentas que
tengan velocidad de viento mayor a 50 nudos, y luego ordénalas de mayor
a menor velocidad de viento. Sugerencia: Usa filter(col \textgreater{}
50) para seleccionar las tormentas fuertes, y arrange(desc(col)) para
ordenar los resultados de mayor a
menor.}\label{q4-filtra-el-dataset-para-quedarte-solo-con-tormentas-que-tengan-velocidad-de-viento-mayor-a-50-nudos-y-luego-orduxe9nalas-de-mayor-a-menor-velocidad-de-viento.-sugerencia-usa-filtercol-50-para-seleccionar-las-tormentas-fuertes-y-arrangedesccol-para-ordenar-los-resultados-de-mayor-a-menor.}

\begin{Shaded}
\begin{Highlighting}[]
\NormalTok{storms }\SpecialCharTok{\%\textgreater{}\%}
  \FunctionTok{filter}\NormalTok{(wind }\SpecialCharTok{\textgreater{}} \DecValTok{50}\NormalTok{) }\SpecialCharTok{\%\textgreater{}\%}
  \FunctionTok{arrange}\NormalTok{(}\FunctionTok{desc}\NormalTok{(wind))}
\end{Highlighting}
\end{Shaded}

\begin{verbatim}
## # A tibble: 7,406 x 13
##    name     year month   day  hour   lat  long status    category  wind pressure
##    <chr>   <dbl> <dbl> <int> <dbl> <dbl> <dbl> <fct>        <dbl> <int>    <int>
##  1 Allen    1980     8     7    18  21.8 -86.4 hurricane        5   165      899
##  2 Gilbert  1988     9    14     0  19.7 -83.8 hurricane        5   160      888
##  3 Wilma    2005    10    19    12  17.3 -82.8 hurricane        5   160      882
##  4 Dorian   2019     9     1    16  26.5 -77   hurricane        5   160      910
##  5 Dorian   2019     9     1    18  26.5 -77.1 hurricane        5   160      910
##  6 Allen    1980     8     5    12  15.9 -70.5 hurricane        5   155      932
##  7 Allen    1980     8     7    12  21   -84.8 hurricane        5   155      910
##  8 Allen    1980     8     8     0  22.2 -87.9 hurricane        5   155      920
##  9 Allen    1980     8     9     6  25   -94.2 hurricane        5   155      909
## 10 Gilbert  1988     9    14     6  19.9 -85.3 hurricane        5   155      889
## # i 7,396 more rows
## # i 2 more variables: tropicalstorm_force_diameter <int>,
## #   hurricane_force_diameter <int>
\end{verbatim}

\paragraph{Q5. Crea un boxplot que muestre la distribución del número de
opalinus bajo diferentes condiciones de luz (Sun vs.~Shade) utilizando
ggplot2. Sugerencia: Usa ggplot(), aes() y
geom\_boxplot()}\label{q5.-crea-un-boxplot-que-muestre-la-distribuciuxf3n-del-nuxfamero-de-opalinus-bajo-diferentes-condiciones-de-luz-sun-vs.-shade-utilizando-ggplot2.-sugerencia-usa-ggplot-aes-y-geom_boxplot}

\end{document}
